%-------------------------
% Resume in Latex
% Author : Jose Sanchez Gonzalez (Modified from Jake Gutierrez's Template)
% License : MIT
%------------------------
%
% PROJECT SWAP GUIDE ------------------------------------------------------
% Four projects fill the page. All six are written below; the two that are
% not currently in use sit inside \begin{comment} ... \end{comment} blocks
% at the bottom of the Projects section. To swap: move the block you want
% out of the comment and move the one you're benching into it.
%
%   ACTIVE (default, generalist AI/ML role):
%     1. Vision-Based Metric Depth Estimation   -- the capstone. Never cut.
%     2. Agentic Soccer Prediction over MCP     -- agents, calibration, MCP
%     3. DoM Decoder-Only Transformer           -- from-scratch transformer
%     4. MetaNaviT                              -- RAG / backend infra
%
%   BENCHED (swap in per role):
%     5. Metric Depth to Volumetric Twin        -- swap in for CV / 3D / robotics
%                                                  roles (replace #4 MetaNaviT)
%     6. 3D Point Cloud Classification          -- swap in for LiDAR / point-cloud
%                                                  roles (replace #2 or #4)
%
%   Quick presets:
%     CV / Robotics / Perception ....... 1, 5, 2, 6   (drop MetaNaviT, DoM)
%     LLM / NLP / Applied Scientist .... 1, 3, 4, 2
%     Agents / LLM Infra / Platform .... 2, 4, 3, 1
%     Generalist / New Grad SWE-ML ..... 1, 2, 3, 4   (default below)
%
% IF IT SPILLS ONTO PAGE 2, cut in this order (each is one line back):
%   1. Depth bullet 3 (the latency profile)
%   2. The "Math and Computer Science Tutor" entry
%   3. The ICPC / Jujitsu line
%   4. Trim the Summary to two sentences
% -------------------------------------------------------------------------

\documentclass[letterpaper,10pt]{article}

\usepackage{latexsym}
\usepackage[letterpaper,top=0.45in,bottom=0.45in,left=0.45in,right=0.45in]{geometry}
\usepackage{titlesec}
\usepackage[usenames,dvipsnames]{color}
\usepackage{enumitem}
\usepackage[hidelinks]{hyperref}
\usepackage{fancyhdr}
\usepackage[english]{babel}
\usepackage{tabularx}
% provides the comment environment used to bench unused projects below
\usepackage{verbatim}
\input{glyphtounicode}

\pagestyle{fancy}
\fancyhf{}
\fancyfoot{}
\renewcommand{\headrulewidth}{0pt}
\renewcommand{\footrulewidth}{0pt}

\urlstyle{same}
\raggedbottom
\raggedright
\setlength{\tabcolsep}{0in}

\titleformat{\section}{
  \vspace{-11pt}\scshape\raggedright\large\bfseries
}{}{0em}{}[\color{black}\titlerule \vspace{-8pt}]

\pdfgentounicode=1

\newcommand{\resumeItem}[1]{
  \item\small{#1 \vspace{-1pt}}
}

\newcommand{\resumeSubheading}[4]{
  \vspace{-2pt}\item
    \begin{tabular*}{\textwidth}[t]{l@{\extracolsep{\fill}}r}
      \textbf{#1} & \textbf{\small #2} \\
      \textit{\small#3} & \textit{\small #4} \\
    \end{tabular*}\vspace{-6pt}
}

\newcommand{\resumeProjectHeading}[2]{
    \item
    \begin{tabular*}{\textwidth}{l@{\extracolsep{\fill}}r}
      \small#1 & \textbf{\small #2}\\
    \end{tabular*}\vspace{-5pt}
}

\newcommand{\resumeSubHeadingListStart}{\begin{itemize}[leftmargin=0.0in, label={}, itemsep=1pt]}
\newcommand{\resumeSubHeadingListEnd}{\end{itemize}\vspace{-1pt}}
\newcommand{\resumeItemListStart}{\begin{itemize}[leftmargin=0.15in, itemsep=1pt, topsep=2pt, parsep=0pt]}
\newcommand{\resumeItemListEnd}{\end{itemize}\vspace{-2pt}}

\begin{document}

%----------HEADING----------
\begin{center}
    {\Huge \scshape Jose Sanchez Gonzalez} \\ \vspace{3pt}
    \small AI/ML Engineer $|$ Deep Learning, Computer Vision, LLM Systems \\ \vspace{3pt}
    \small (503) 660-9763 ~$|$~
    \href{mailto:joseszgz2021@gmail.com}{joseszgz2021@gmail.com} ~$|$~
    \href{https://linkedin.com/in/jose-j-sanchez-gonzalez-84a800257/}{linkedin.com/in/jose-j-sanchez-gonzalez} ~$|$~
    \href{https://jose-sanchez-portfolio-com.vercel.app/}{jose-sanchez-portfolio.vercel.app}
\end{center}
\vspace{-12pt}

%-----------SUMMARY-----------
\section{Summary}
\resumeSubHeadingListStart
  \item \small{AI/ML Engineer with an M.S. in Artificial Intelligence. My work centers on the systematic development and benchmarking of deep learning models, emphasizing reproducibility and quantitative evaluation. I work primarily in PyTorch, addressing challenges in computer vision, 3D perception, and large language model systems. Recent projects include metric depth estimation for orchard robotics, a transformer built from first principles, calibrated agentic forecasting served over MCP, and containerized retrieval-augmented generation services.}
\resumeSubHeadingListEnd

%-----------EXPERIENCE-----------
\section{Experience}
\resumeSubHeadingListStart
  \resumeSubheading
    {Oregon State University}{Corvallis, OR}
    {Graduate Teaching Assistant}{Sep 2025 -- Jun 2026}
    \resumeItemListStart
      \resumeItem{Supported Operating Systems I (C++) and Computational Methods (COMSOL), covering concurrency, memory management, file systems, numerical methods, and simulation workflows.}
      \resumeItem{Led 2-hour studio sessions for 25+ students, held office hours, and provided programming, debugging, and modeling support.}
    \resumeItemListEnd
  \resumeSubheading
    {Oregon State University}{Corvallis, OR}
    {Math and Computer Science Tutor}{Nov 2022 -- Jun 2025}
    \resumeItemListStart
      \resumeItem{Tutored 50+ students in math, programming, and CS fundamentals, adapting technical explanations to different backgrounds and skill levels.}
    \resumeItemListEnd
\resumeSubHeadingListEnd

%-----------SKILLS-----------
\section{Skills}
\resumeSubHeadingListStart
  \item \small{
    \textbf{Languages:} Python, C++, C, SQL, Java, Kotlin \\
    \textbf{ML \& Deep Learning:} PyTorch, scikit-learn, XGBoost, NumPy, Pandas, CNNs, Transformers, Attention, LSTMs/GRUs, Autoencoders, Reinforcement Learning \\
    \textbf{Computer Vision \& 3D:} DINOv2, Vision Transformers, OpenCV, PointNet, RGB-D \& Stereo Depth Estimation, Camera Calibration, Blender \\
    \textbf{LLM \& Agents:} Retrieval-Augmented Generation (RAG), LangGraph, Model Context Protocol (MCP), LlamaIndex, Ollama, pgvector, Vector Search, Constrained Decoding \\
    \textbf{Infrastructure:} Docker, FastAPI, PostgreSQL, MySQL, MongoDB, Redis, Linux, Git, ROS2, NVIDIA Isaac Sim, Slurm/HPC
  }
\resumeSubHeadingListEnd

%-----------EDUCATION-----------
\section{Education}
\resumeSubHeadingListStart
  \resumeSubheading
    {Oregon State University}{Corvallis, OR}
    {M.S. in Artificial Intelligence}{Graduated Jun 2026}
  \resumeSubheading
    {Oregon State University}{Corvallis, OR}
    {B.S. in Computer Science (Applied Route, Data Science), Mathematics Minor}{Graduated Jun 2025}
\resumeSubHeadingListEnd

%-----------PROJECTS-----------
% See the PROJECT SWAP GUIDE at the top of this file.
\section{Projects}
\resumeSubHeadingListStart

  % ---- 1. CAPSTONE: metric depth for robotic pruning ---------------------
  \resumeProjectHeading
    {\href{https://jose-sanchez-portfolio-com.vercel.app/projects/depth-estimation-robotic-pruning}{\textbf{Vision-Based Metric Depth Estimation for Robotic Pruning}} $|$ \emph{Python, PyTorch, DINOv2, CNNs, Blender, OpenCV}}{}
    \resumeItemListStart
      \resumeItem{Built a synthetic-orchard pipeline in Blender rendering RGB, per-pixel ground-truth depth, masks, and full camera intrinsics/extrinsics, then benchmarked monocular, scale-calibrated, fine-tuned, and multi-view RGB-D models under one protocol: 10 trees, 5 seeds, fixed frame budget, one design change at a time.}
      \resumeItem{DINOv2 RGB+D multi-view refinement over 3 stereo pairs cut validation RMSE from 0.0550 m to $0.0445 \pm 0.0057$ m; a streaming least-squares scale/shift fit aligned predictions to metric depth in one constant-memory pass over 24k+ images.}
      \resumeItem{Profiled deployment latency alongside accuracy: 484.7 ms per 6-view group (80.8 ms amortized per depth map) vs. 333.6 ms per monocular frame, making the speed/precision trade-off explicit.}
    \resumeItemListEnd

  % ---- 2. Agentic forecasting over MCP -----------------------------------
  \resumeProjectHeading
    {\href{https://jose-sanchez-portfolio-com.vercel.app/projects/agentic-soccer-mcp}{\textbf{Agentic Soccer Prediction over MCP}} $|$ \emph{Python, LangGraph, MCP, XGBoost, Conformal Prediction, FastAPI, Docker}}{}
    \resumeItemListStart
      \resumeItem{Built an offline forecasting stack---leakage-guarded features into XGBoost with isotonic calibration and split-conformal prediction sets, reconciled against a Dixon-Coles score grid---and served it through a FastAPI gateway into a LangGraph orchestrator driving three MCP servers (odds, news, inference), with a human-in-the-loop interrupt before any staking suggestion.}
      \resumeItem{On 102 World Cup 2026 matches, trained strictly pre-tournament: 0.9013 log loss and 62.7\% accuracy vs. 47.1\% for the prior, with split-conformal coverage at 0.951 against a 0.90 target.}
    \resumeItemListEnd

  % ---- 3. DoM: decoder-only transformer ----------------------------------
  \resumeProjectHeading
    {\href{https://jose-sanchez-portfolio-com.vercel.app/projects/dom-math-transformer}{\textbf{DoM: Decoder-Only Transformer for Math Word Problem Translation}} $|$ \emph{Python, PyTorch, NLP, Transformers}}{}
    \resumeItemListStart
      \resumeItem{Collaboratively implemented a decoder-only sequence-to-sequence transformer in PyTorch with a custom tokenizer, vocabulary, and dataloader, translating $\sim$30,000 MathQA word problems into executable computational graphs.}
      \resumeItem{Designed a constrained decoding method that guarantees valid directed acyclic graph construction, reaching 77.45\% exact-match accuracy and 84.78\% BLEU on the test set.}
    \resumeItemListEnd

  % ---- 4. MetaNaviT: RAG file organizer ----------------------------------
  \resumeProjectHeading
    {\href{https://jose-sanchez-portfolio-com.vercel.app/projects/metanavit}{\textbf{MetaNaviT: LLM-Powered File Organizer}} $|$ \emph{Python, FastAPI, PostgreSQL, pgvector, LlamaIndex, Ollama, Docker}}{}
    \resumeItemListStart
      \resumeItem{Built a containerized retrieval-augmented generation service that indexes local, cloud, and web sources, using pgvector and locally hosted LLMs for semantic search, metadata extraction, hierarchical indexing, and automated file structuring.}
    \resumeItemListEnd

% ==========================================================================
% BENCHED PROJECTS --- move a block out of \begin{comment}...\end{comment}
% to activate it, and comment out whichever active project it replaces.
% ==========================================================================

\begin{comment}

  % ---- 5. Metric Depth to Volumetric Twin --------------------------------
  % Swap in for computer vision / 3D / sports-tech / research roles.
  \resumeProjectHeading
    {\href{https://jose-sanchez-portfolio-com.vercel.app/projects/metric-depth-volumetric-twin}{\textbf{Metric Depth to Volumetric Twin}} $|$ \emph{Python, PyTorch, YOLO, Camera Calibration, Triangulation}}{}
    \resumeItemListStart
      \resumeItem{Reproduced the SoccerNet-v3D (CVPR CVSports 2025) detection and monocular 3D localization baselines on the paper's own frozen splits---within 0.003 AP of the published figure, with reprojection matched to a median of 0.002 px across 4,051 frames.}
      \resumeItem{Replaced blur-corrupted size-prior depth with a gravity-constrained ballistic trajectory fit over a 360 ms window, cutting median depth error from 3.26 m to 2.82 m and returning 5.4$\times$ more estimates inside 2 m.}
    \resumeItemListEnd

  % ---- 6. 3D Point Cloud Classification ----------------------------------
  % Swap in for LiDAR / point-cloud / autonomous-driving roles.
  \resumeProjectHeading
    {\href{https://jose-sanchez-portfolio-com.vercel.app/projects/point-cloud-classification}{\textbf{3D Point Cloud Classification}} $|$ \emph{Python, PyTorch, LiDAR, Deep Learning}}{}
    \resumeItemListStart
      \resumeItem{Implemented models inspired by PointNet and PointConvFormer, incorporating attention over local neighborhoods to classify 3D objects on the ModelNet and ShapeNet benchmarks; attention variants improved over the PointNet-style baseline.}
    \resumeItemListEnd

\end{comment}

\resumeSubHeadingListEnd

%-----------LEADERSHIP-----------
\section{Leadership \& Activities}
\resumeItemListStart
  \resumeItem{\textbf{Berkeley Humanoid Lite Build} (in progress): Independently constructing an open-source humanoid robot, focusing on mechanical assembly, soldering, and CAD design. Integration of ROS2 and NVIDIA Isaac Sim aims to enable motion control and reinforcement learning.}
  \resumeItem{\textbf{ICPC Programming Contest Club} and \textbf{Jujitsu Club}: Competitive programming practice; regular training.}
\resumeItemListEnd

\end{document}
